\section{3. Pythagoras' sætning}

\kompetence{Regning} \kompetence{Tankegang}

\begin{figure}[ht]
\centering
\begin{tikzpicture}[>=Stealth,scale=2.2]
  % Retvinklet trekant
  \coordinate (A) at (0,0);
  \coordinate (B) at (2.4,0);
  \coordinate (C) at (0,1.8);

  % Trekant
  \draw[very thick] (A) -- (B) node[midway,below] {$a$} -- (C) node[midway,right] {$c$} -- cycle;

  % Højre vinkel-markør
  \draw (0,0.18) -- (0.18,0.18) -- (0.18,0);

  % Vinkel-hypotenuse navn
  \node[anchor=south west] at (0.6,0.35) {$\theta$};

  % Katet-navne
  \node[below left] at (0.6,-0.05) {$b$};

  % Areal-tekster
  \node[anchor=south] at (1.2,0) {$\text{Areal}=a\cdot b/2$};
\end{tikzpicture}
\qquad
\begin{tikzpicture}[>=Stealth,scale=2.2]
  % Klassisk Pythagoras-bevis: stor kvadrat (a+b)² med fire trekanter
  % Trekant-hjørner
  \coordinate (A) at (0,0);
  \coordinate (B) at (4.2,0);
  \coordinate (C) at (4.2,3.2);
  \coordinate (D) at (0,3.2);
  % Trekanthjørner indefra
  \coordinate (T1) at (1.2,0);
  \coordinate (T2) at (4.2,1.2);
  \coordinate (T3) at (3.0,3.2);
  \coordinate (T4) at (0,2.0);
  % Diagonal kvadrat hjørner
  \coordinate (H1) at (1.2,1.2);
  \coordinate (H2) at (3.0,1.2);
  \coordinate (H3) at (3.0,2.0);
  \coordinate (H4) at (1.2,2.0);

  % Store kvadrat yderside
  \draw[thick] (A) rectangle (C);
  % Fire trekanter
  \draw[fill=gxu-accent!20] (A) -- (T1) -- (H1) -- (H4) -- cycle;
  \draw[fill=gxu-accent!20] (T1) -- (B) -- (H2) -- (H1) -- cycle;
  \draw[fill=gxu-accent!20] (H2) -- (C) -- (T3) -- (H3) -- cycle;
  \draw[fill=gxu-accent!20] (H4) -- (H3) -- (T4) -- (D) -- cycle;
  % Indre kvadrat (hypotenusen c)
  \draw[very thick,gxu-purple] (H1) -- (H2) -- (H3) -- (H4) -- cycle;
  % Areal-tekst
  \node[anchor=north east,font=\sffamily\small] at (H1) {$a^{2}$};
  \node[anchor=north west,font=\sffamily\small] at (H3) {$b^{2}$};
  \node[anchor=center,font=\sffamily\small\color{gxu-purple}] at (H2 |- H3) {$c^{2}$};
\end{tikzpicture}
\caption{Venstte: retvinklet trekant med kateter $a$ og $b$ samt hypotenuse $c$. Højre: samme trekant indlagt i et kvadrat, der viser at $a^{2}+b^{2}=c^{2}$.}
\label{fig:pythagoras}
\end{figure}

For en retvinklet trekant gælder:

\eq{3.1}{a^{2}+b^{2}=c^{2}}

hvor $c$ er hypotenusen (den længste side, modsat den rette vinkel).

\begin{prompt}
En retvinklet trekant har kateter $a=3$ og $b=4$. Find hypotenusen $c$.
\end{prompt}
\begin{solutionc}
$c^{2}=3^{2}+4^{2}=9+16=25$, så $c=5$.
\end{solutionc}
